% !TEX root = ../main.tex
\section{Discussion and Conclusion}
\label{sec:discussion}

This chapter answers the research question from the results of
Section~\ref{sec:results}, discusses the materials that depart from the average, and
closes with the limitations of the study, possible future work and a conclusion.

% ============================================================================
\subsection{Answering the Research Question}
\label{sec:discussion-answer}

The research question has two parts: at which O2I ratio the complex background of the
learning factory begins to dominate the classification decision, and whether
explainability methods can detect this shift. Section~\ref{sec:methodology-rq} sets out
how the results are read for each part, and the answer below follows that reading.

\paragraph{Classification}
Averaged over the 15 materials, the test F1 falls from 0.813 at an O2I ratio of 90\% to
0.720 at 10\% (Tab.~\ref{tab:ovr-full-results}), and it falls in 14 of the 15
materials. Most of this fall lies below 50\%. From 90\% to 50\% the mean changes by
$-0.017$, and the paired test shows no difference ($p = 0.25$), whereas from 50\% to
10\% it falls clearly, by 0.076 ($p = 0.001$). After 200 epochs the validation F1
still falls with the ratio, from 0.884 at 90\% to 0.841 at 10\%, against 0.854 to
0.794 after 100 epochs, so the training length explains a small part of the decline,
not the decline itself.

\paragraph{Two caveats}
The mean ROC-AUC stays between 0.970 and 0.982 at every ratio, so the models still rank
most positive crops above most negative ones while the F1 at the threshold $p_1 = 0.5$
falls. The decline of the F1 also does not isolate the background. With a fixed input
of $256 \times 256$ pixels, the side of a square with the area of the box shrinks from
about 243 pixels at 90\% to about 81 pixels at 10\%, i.e., from about 6.4 to about 2.1
times the 38-pixel receptive field of the fourth convolutional layer (paragraph
\emph{What changes with the ratio} in Section~\ref{sec:methodology-crops}). The F1
decline fits a growing influence of the background, but it cannot separate that
influence from the smaller scale of the workpiece.

\paragraph{Where the evidence lies}
The Grad-CAM centre measure falls at every step, from 0.737 at 90\% to 0.407 at 10\%
(Tab.~\ref{tab:centerfrac}). At these mean values, Eq.~\eqref{eq:odds} gives
$m_{\mathrm{in}}/m_{\mathrm{out}} = 2.80$ and 0.69. Since every reference value rises as
the ratio falls (Tab.~\ref{tab:reference}), the fall can come neither from the geometry
of the square nor from the resolution of Grad-CAM, and it is read as evidence moving
from the workpiece into the background. Its largest step, from 30\% to 10\%, lies in
the same interval as the largest step of the F1. This is a description of the two
curves, not a tested result (Section~\ref{sec:methodology-tests}).

Under the rule of Section~\ref{sec:methodology-rq}, the background outweighs the
workpiece only where $c$ lies clearly below 0.5, and this holds at no ratio for either
method. At 10\%, Grad-CAM is below 0.5 on average and in 9 of the 15 materials, but not
clearly ($t(14) = -1.36$), and the same holds without the flat maps (0.363,
$t(14) = -1.64$). The criterion is strict by design. At 10\% the region outside the
reference square has about nine times its area, so the background holds more of the
total importance whenever $c$ is below about 0.9. By that total, even occlusion would
show the background dominating at 10\% on average, since its mean centre measure of
0.884 lies below 0.9 although it strongly favours the square per pixel. The total depends on the area of
the square and would change with the ratio even for a flat map, which is why the
per-pixel measure is used (Section~\ref{sec:methodology-rq}). The answer to the first
part of the research question is therefore a gradual shift of the evidence towards the
background over the five ratios, without a ratio at which the background clearly
dominates, although the classification suffers mainly below 50\%.

\paragraph{Detection of the shift}
By the detection rule of Section~\ref{sec:methodology-rq}, Grad-CAM detects the shift:
its centre measure changes by $-0.330$ from 90\% to 10\% ($t(14) = -4.17$) while the
test F1 changes by $-0.093$ ($t(14) = -4.76$). Occlusion sensitivity does not, since
its centre measure rises from 0.728 to 0.884. The rise does not show that the occlusion
evidence stays on the workpiece. The occlusion reference value rises faster, from 0.569
to 0.956, so the reference gap of occlusion falls from $+0.159$ at 90\% to about zero or
below from 50\% on. This fall is at most an indirect sign of a shift, since the
reference gaps are not tested as trends (Section~\ref{sec:methodology-refgap}).

The two methods agree at 90\% and then move apart. The method difference is clear at
10\% (0.477) and as a slope over the five ratios, and at 10\% the difference beyond the
reference is 0.509, so the different resolutions of the two methods do not produce it
(Section~\ref{sec:results-reference}). Grad-CAM thus detects the shift and occlusion
sensitivity, the method specified for this thesis, does not, and whether a shift of the
evidence is visible depends on the method used to look for it. None of the five checks of
Section~\ref{sec:results-robustness} changes this result. The agreement at 90\% does
not hold under all of them: it holds only for positive drops with the grey fill, while
with min-max normalisation or a blurred patch occlusion lies clearly below Grad-CAM at
90\%.

\paragraph{Why the two methods answer differently}
Section~\ref{sec:background-comparison} gives the reason why the two methods can
disagree. If the background holds many similar cues, like a room full of clocks,
covering one of them with a patch changes the score little, while Grad-CAM marks all of
them at once. The results add two points to this. At 10\% the workpiece fits into a
reference square of about 81 pixels, and one $32 \times 32$ patch covers about 40\% of
its width, so a few patches remove most of the workpiece and the drops concentrate
there. Together with the rising reference value, this would explain why the occlusion
centre measure even rises as the ratio falls. The fill of the patch, in contrast, cannot
explain the average divergence, since the divergence appears with the grey, the black
and the blurred patch alike (Section~\ref{sec:results-fill}); the fill matters for
single workpieces (Section~\ref{sec:discussion-material}). The redundancy explanation
remains a hypothesis. The thesis measures that the methods diverge, not why. The toy of
Section~\ref{sec:background-comparison} shows only the ingredients, and neither map is
ground truth for what the model uses. A direct test would replace the whole background
at once instead of one patch (Section~\ref{sec:discussion-future}).

\paragraph{On average, not per material}
The detection rule compares means over the 15 materials. It shows that the mean F1 and
the mean Grad-CAM centre measure fall together, not that the materials whose
classification suffers most are those whose evidence moves most, and per material no
such link is shown. Over the 15 materials, the change from 90\% to 10\% of the test F1
(Tab.~\ref{tab:ovr-full-results}) and that of the per-model Grad-CAM centre measures
behind Tab.~\ref{tab:centerfrac} give a Pearson correlation of $r = -0.20$
($p = 0.47$) and a Spearman rank correlation of $\rho = -0.14$ ($p = 0.63$). Steel
loses most F1, 0.230, while its Grad-CAM centre measure rises by 0.148, aluminium loses
0.224 while its centre measure falls by only 0.074, and blue AG has the largest
Grad-CAM fall, 0.827, with almost no loss of F1 (0.028). Grad-CAM thus shows that the
evidence shifts towards the background as the background grows, not that this shift is
why a given material is misclassified (Section~\ref{sec:discussion-material}).

\paragraph{Relation to earlier work}
The gradual decline of the F1, without a fall to chance, fits the results of Pawlowski
et al.~\cite{pawlowski2019needles} for ratio ranges that cover those studied here
(Section~\ref{sec:related-work-o2i}). Unlike their saliency maps, which fail where the
classifier fails, the Grad-CAM evidence here moves away from the workpiece while the
classification declines only moderately and the models still rank the crops well
(mean ROC-AUC of at least 0.970). Xiao et al.~\cite{xiao2021noise} showed that the
background can carry evidence for a class, so evidence in the background is not by
itself an error. Their further finding, that more accurate models exploit the
background more, compares different architectures on the same images, whereas here one
architecture is trained on crops of different ratios.

% ============================================================================
\subsection{Material-Dependent Effects}
\label{sec:discussion-material}

[6.2 to follow]

% ============================================================================
\subsection{Limitations}
\label{sec:discussion-limitations}

[6.3 to follow]

% ============================================================================
\subsection{Future Work}
\label{sec:discussion-future}

[6.4 to follow]

% ============================================================================
\subsection{Conclusion}
\label{sec:discussion-conclusion}

[6.5 to follow]
